\documentclass[11pt]{article}
\usepackage[margin=1in]{geometry}
\usepackage{amsmath,amssymb}
\usepackage{booktabs}
\usepackage{graphicx}
\usepackage{algorithm}
\usepackage{algpseudocode}
\usepackage{hyperref}
\usepackage{xcolor}

\title{RLE-Based Matrix--Vector Multiplication\\from Threading Instructions}
\author{Technical Note}
\date{March 2026}

\begin{document}
\maketitle

\begin{abstract}
We describe a run-length encoded (RLE) approach to genotype-matrix--vector
multiplication from the \textsc{Threads} ARG representation.  The key
observation is that the implicit genotype matrix encoded by threading
instructions can be decomposed into per-sample lists of contiguous
``1-runs'' during a single preparation pass, after which both left- and
right-multiplication reduce to prefix-sum lookups over these runs.  On
simulated data with 10{,}000 diploid individuals, the RLE multiply is
\textbf{15--17$\times$ faster} than the existing tree-propagation multiply
and \textbf{14--60$\times$ faster} than dense multiplication, while
requiring comparable preparation time and modest additional memory.  The
method is embarrassingly parallel and admits efficient batched
multiplication over $k$ vectors simultaneously.
\end{abstract}

\section{Background}

The \textsc{Threads} software~\cite{gunnarsson2025} infers an Ancestral
Recombination Graph (ARG) by threading each haplotype through a set of
reference haplotypes.  The result is a set of \emph{threading instructions}:
for each of $n$ haplotype samples, a sequence of segments
$\{(\text{start}_k, \text{target}_k, \text{tmrca}_k)\}$ and a list of
mismatch site indices.  These instructions implicitly encode the $m \times n$
binary genotype matrix $\mathbf{G}$, where $G_{ji}$ is the allele of sample
$i$ at site $j$.

Many downstream applications---RHE~\cite{pazokitoroudi2020}, SCORE, and
SCOPE---require repeated matrix--vector products $\mathbf{G}^\top
\mathbf{x}$ (right multiply) and $\mathbf{G}\mathbf{x}$ (left multiply).
Three existing approaches are available:

\begin{enumerate}
\item \textbf{Dense multiply}: Materialize $\mathbf{G}$ in memory
($O(nm)$ space), then use standard matrix--vector products.  Per-call cost:
$O(nm)$.
\item \textbf{Tree-propagation multiply}: Exploit the tree structure of
threading instructions.  A one-time $O(nm)$ preparation step (with $O(m
\cdot d)$ peak memory, where $d$ is the tree depth) precomputes cumulative
interval sums and mismatch corrections.  Per-call cost: $O(n \cdot
n_\text{intervals} + n_\text{mismatches})$, where $n_\text{intervals}$ is
the number of unique segment boundaries across all samples.
\item \textbf{GRG multiply}~\cite{dehaas2024}: Build a Genotype
Representation Graph (DAG) and propagate partial sums.  Per-call cost:
$O(|\text{edges}| + m) \approx O(n \log n + m)$.
\end{enumerate}

The tree-propagation approach has a fundamental limitation: the cumulative
sum for each sample depends on its target's cumulative sum, creating a
sequential dependency chain.  This prevents parallelization and results in
$O(n \cdot n_\text{intervals})$ work even though only $O(n_\text{segments}
+ n_\text{mismatches})$ information determines the genotype matrix.

\section{Method}

\subsection{Run-Length Encoding of Genotypes}

For each sample $i$, the binary genotype vector $\mathbf{g}_i \in
\{0,1\}^m$ consists of contiguous runs of 0s and 1s.  A \emph{1-run} is a
maximal interval $[a, b)$ where $G_{ji} = 1$ for all $a \le j < b$.  We
store only the 1-runs:
\[
\mathcal{R}_i = \{(a_1, b_1), (a_2, b_2), \ldots, (a_{R_i}, b_{R_i})\}
\]
where $R_i = |\mathcal{R}_i|$ is the number of 1-runs for sample $i$.

The total number of runs across all samples is
$R_\text{total} = \sum_i R_i$.  In practice, $R_i$ is determined by the
number of transitions in $\mathbf{g}_i$, which depends on the number of
segments and mismatches for sample $i$ and its ancestral lineage.

\subsection{Preparation Phase}

The 1-run lists are constructed during a single pass over samples in
threading order (sample 0 first, then 1, 2, \ldots, $n{-}1$), using a
reference-counted genotype cache identical to the existing tree-propagation
preparation.

For sample 0 (the reference haplotype), the genotype is determined directly
from its mismatch list.  For sample $i > 0$, each segment copies the
genotype from its target (already computed and cached) and flips at mismatch
sites.  After reconstructing the genotype buffer, a single linear scan
extracts the 1-runs.

\begin{algorithm}[h]
\caption{RLE Preparation}
\begin{algorithmic}[1]
\For{$i = 0$ to $n-1$}
    \State Reconstruct $\mathbf{g}_i$ from target genotypes + mismatches
    \State Scan $\mathbf{g}_i$ to extract 1-runs $\mathcal{R}_i$
    \State Store $\mathcal{R}_i$ in flat arrays with offset index
    \If{sample $i$ is referenced by future samples}
        \State Cache $\mathbf{g}_i$
    \EndIf
    \State Release unreferenced target caches
\EndFor
\end{algorithmic}
\end{algorithm}

The preparation cost is $O(nm)$ time with $O(m \cdot d_\text{max})$ peak
memory for the genotype cache, where $d_\text{max}$ is the maximum number
of simultaneously active targets (bounded by the tree depth).  The stored
1-runs require $O(R_\text{total})$ memory.

\subsection{Right Multiply: $\mathbf{G}^\top \mathbf{x}$}

Given a vector $\mathbf{x} \in \mathbb{R}^m$, compute $\mathbf{y} =
\mathbf{G}^\top \mathbf{x}$, where $y_i = \sum_j G_{ji} \, x_j$.

\begin{enumerate}
\item Compute prefix sums: $P_0 = 0$, $P_s = \sum_{j=0}^{s-1} x_j$ for
$s = 1, \ldots, m$.
\item For each sample $i$:
\[
y_i = \sum_{(a,b) \in \mathcal{R}_i} (P_b - P_a)
\]
\end{enumerate}

Cost: $O(m)$ for prefix sums $+$ $O(R_\text{total})$ for the per-sample
summation.  The per-sample loop is \textbf{embarrassingly parallel} ---
each sample's computation is independent.

\subsection{Left Multiply: $\mathbf{G} \mathbf{x}$}

Given $\mathbf{x} \in \mathbb{R}^n$, compute $\mathbf{y} = \mathbf{G}
\mathbf{x}$, where $y_j = \sum_i G_{ji} \, x_i$.

\begin{enumerate}
\item Initialize a difference array $D[0 \ldots m] = 0$.
\item For each sample $i$, for each $(a,b) \in \mathcal{R}_i$:
$D[a] \mathrel{+}= x_i$, \; $D[b] \mathrel{-}= x_i$.
\item Compute prefix sums: $y_j = \sum_{s=0}^{j} D[s]$.
\end{enumerate}

Cost: $O(R_\text{total})$ for the difference updates $+$ $O(m)$ for the
prefix sum.  With OpenMP, the per-sample loop uses thread-local difference
arrays that are reduced before the final prefix sum.

\subsection{Batch Multiply}

For $k$ vectors simultaneously (input: $m \times k$ or $n \times k$
row-major matrices), the prefix sums and difference arrays are extended to
width $k$.  The inner loops over $k$ elements auto-vectorize via SIMD,
amortizing the run-traversal overhead.

\subsection{Simplification}

As a complementary optimization, we also implement
\texttt{simplify\_for\_multiply()}, which strips coalescence times
(\texttt{tmrca} values, unused by any multiply algorithm) and merges
consecutive segments sharing the same target.  This yields $\sim$30\%
smaller \texttt{.threads} files with no effect on the genotype matrix.
However, segment merging is minimal (0.2--2\% reduction) because
consecutive same-target segments are rare in real ARGs.

\section{Results}

We benchmark on simulated data (msprime, $N_e = 10{,}000$, 2\,Mbp, seed
42) at sample sizes from 50 to 10{,}000 diploid individuals.  We compare
five methods: RLE (this work, with OpenMP), tree-propagation multiply
(existing \textsc{Threads} method), dense materialized multiply,
GRG~\cite{dehaas2024}, and RePair grammar-compressed
multiply~\cite{tosoni2022}.  All timings are median of 5 runs on an Apple
M-series processor.

\subsection{Right Multiply $\mathbf{G}^\top \mathbf{x}$ (ms)}

\begin{table}[h]
\centering
\begin{tabular}{r rrrrr}
\toprule
$n_\text{dip}$ & RLE & Tree & Dense & GRG & RePair \\
\midrule
50    & \textbf{0.10}  & 0.46  & 0.11  & 0.30  & ---   \\
100   & \textbf{0.13}  & 0.66  & 0.23  & 0.40  & 0.32  \\
250   & \textbf{0.17}  & 0.98  & 0.55  & 0.55  & ---   \\
500   & \textbf{0.24}  & 1.38  & 1.15  & 0.70  & ---   \\
1000  & \textbf{0.44}  & 3.33  & 2.50  & 0.94  & 2.40  \\
2000  & \textbf{0.72}  & 7.04  & 5.34  & 1.12  & ---   \\
5000  & 1.65  & 21.68 & 14.63 & \textbf{1.55}  & ---   \\
10000 & 2.90  & 46.61 & 44.45 & \textbf{1.90}  & 25.86 \\
\bottomrule
\end{tabular}
\caption{Right multiply times (ms).  RLE is fastest up to $\sim$5K diploid;
GRG overtakes at larger sizes due to its $O(n\log n)$ DAG traversal.}
\label{tab:right}
\end{table}

\subsection{Left Multiply $\mathbf{G} \mathbf{x}$ (ms)}

\begin{table}[h]
\centering
\begin{tabular}{r rrrrr}
\toprule
$n_\text{dip}$ & RLE & Tree & Dense & GRG & RePair \\
\midrule
50    & \textbf{0.17}  & 0.48  & 0.36  & 0.27  & ---   \\
100   & \textbf{0.19}  & 0.82  & 0.81  & 0.40  & 0.31  \\
250   & \textbf{0.22}  & 1.46  & 2.99  & 0.56  & ---   \\
500   & \textbf{0.35}  & 1.53  & 6.83  & 0.67  & ---   \\
1000  & \textbf{0.52}  & 3.94  & 16.63 & 0.92  & 2.34  \\
2000  & \textbf{0.81}  & 8.43  & 36.20 & 1.13  & ---   \\
5000  & 1.58  & 28.10 & 101.89& \textbf{1.55}  & ---   \\
10000 & 2.55  & 55.00 & 220.30& \textbf{2.00}  & 24.11 \\
\bottomrule
\end{tabular}
\caption{Left multiply times (ms).  Same pattern: RLE dominates up to
$\sim$5K; GRG is comparable or faster at 10K.  Dense is 86$\times$ slower
than RLE at 10K.}
\label{tab:left}
\end{table}

\subsection{Batch Multiply ($k = 10$ vectors, ms)}

\begin{table}[h]
\centering
\begin{tabular}{r rr c rr}
\toprule
& \multicolumn{2}{c}{Right ($k{=}10$)} && \multicolumn{2}{c}{Left ($k{=}10$)} \\
\cmidrule{2-3} \cmidrule{5-6}
$n_\text{dip}$ & RLE & Tree && RLE & Tree \\
\midrule
1000  & 4.5  & 43.6  && ---  & ---  \\
5000  & 12.1 & 209.0 && ---  & ---  \\
10000 & 20.7 & 369.3 && 23.0 & 467.4\\
\bottomrule
\end{tabular}
\caption{Batch multiply ($k = 10$) times (ms).  RLE batch is
10--18$\times$ faster than tree batch.}
\label{tab:batch}
\end{table}

\subsection{Storage}

\begin{table}[h]
\centering
\begin{tabular}{r rrrr}
\toprule
$n_\text{dip}$ & \texttt{.threads} & GRG & RePair & Dense \\
\midrule
100   & \textbf{101\,KB}  & 332\,KB  & 116\,KB  & 1.3\,MB \\
1000  & \textbf{188\,KB}  & 642\,KB  & 1.3\,MB  & 17.5\,MB \\
10000 & \textbf{438\,KB}  & 1.3\,MB  & 13.8\,MB & 219\,MB \\
\bottomrule
\end{tabular}
\caption{On-disk storage. The \texttt{.threads} format is the most compact:
3$\times$ smaller than GRG and 31$\times$ smaller than RePair at 10K
diploid. Simplified \texttt{.threads} (stripping tmrcas) gives an
additional $\sim$30\% reduction.}
\label{tab:size}
\end{table}

\subsection{Scaling Analysis}

RLE is the fastest method up to $\sim$5{,}000 diploid.  Beyond that, GRG's
$O(n\log n)$ DAG traversal becomes advantageous over RLE's $O(R_\text{total})$
scan, as $R_\text{total}$ grows linearly with $n$ while the DAG amortizes
shared ancestry.  At 10K diploid, GRG is 1.5$\times$ faster than RLE for
right multiply and 1.3$\times$ faster for left.

\begin{table}[h]
\centering
\begin{tabular}{r rrrr}
\toprule
$n_\text{dip}$ & RLE/Tree & RLE/Dense & RLE/RePair & RLE/GRG \\
\midrule
100   & 5.1$\times$ & 1.8$\times$ & 2.5$\times$ & 3.1$\times$ \\
1000  & 7.6$\times$ & 5.7$\times$ & 5.5$\times$ & 2.1$\times$ \\
10000 & 16.1$\times$ & 15.3$\times$ & 8.9$\times$ & 0.65$\times$ \\
\bottomrule
\end{tabular}
\caption{Right multiply speedup ratios (RLE time / other method time,
$>$1 means RLE is faster).  RLE beats all methods up to $\sim$5K diploid
and remains competitive with GRG at 10K.}
\label{tab:speedup}
\end{table}

Key observations:
\begin{itemize}
\item Tree multiply is \textbf{inherently sequential} ($O(n \cdot
n_\text{intervals})$ with dependency chain), so RLE's parallelism provides
an increasing advantage: 5$\times$ at 100 dip to 16$\times$ at 10K dip.
\item Dense multiply is memory-bandwidth-bound at $O(nm)$ per call.  At
10K diploid the matrix is 219\,MB, so left multiply (column-major access)
suffers severely.
\item RePair multiply is 9$\times$ slower than RLE and 14$\times$ slower
than GRG at 10K, due to grammar decompression overhead.
\item GRG has the best \emph{asymptotic} scaling but RLE has lower constant
factors, making it faster at moderate sample sizes.
\end{itemize}

\section{Complexity Comparison}

\begin{table}[h]
\centering
\begin{tabular}{l c c c}
\toprule
Method & Prepare & Per-call & Parallel? \\
\midrule
Dense        & $O(nm)$ & $O(nm)$ & Yes \\
Tree         & $O(nm)$ & $O(n \cdot n_i + M)$ & No \\
RLE          & $O(nm)$ & $O(m + R)$ & Yes \\
GRG          & $O(nm)$ & $O(n \log n + m)$ & Partial \\
\bottomrule
\end{tabular}
\caption{Asymptotic complexity.  $n_i$: number of global intervals;
$M$: total mismatches; $R$: total 1-runs.  In practice $R \ll n \cdot
n_i$ and RLE's per-call cost is embarrassingly parallel.}
\label{tab:complexity}
\end{table}

\section{Memory}

The RLE representation stores $R_\text{total}$ run-start and run-end
integers plus $n{+}1$ offset integers.  At 10K diploid with $R_\text{total}
\approx 16$M:
\[
\text{RLE storage} = 2 \times 16\text{M} \times 4\,\text{bytes} \approx 128\,\text{MB}
\]
compared to:
\begin{itemize}
\item Dense: $nm \approx 230$\,MB (int8) or $920$\,MB (int32)
\item Tree: $n \cdot n_i \times 5\,\text{bytes} \approx 175$\,MB
(\texttt{ivl\_seg} + \texttt{carry} + auxiliary)
\end{itemize}

The RLE memory is comparable to tree multiply but enables much faster
per-call execution.

\section{Discussion}

The RLE approach achieves its speedup from two sources:

\textbf{1.\ Algorithmic:} The per-call cost $O(m + R_\text{total})$ avoids
the $O(n \cdot n_\text{intervals})$ cumulative-sum propagation of tree
multiply.  While $R_\text{total}$ is not dramatically smaller than $n
\cdot n_\text{intervals}$ (16M vs 35M at 10K), the operations are simpler
(one subtraction per run vs.\ memcpy + offset adjustment per interval).

\textbf{2.\ Parallelism:} RLE right multiply is embarrassingly parallel
(independent per-sample computation), while tree multiply has inherent
sequential dependencies.  With $p$ cores, the effective cost ratio is
$(R_\text{total}/p) / (n \cdot n_i)$, explaining the increasing speedup
with sample size.

\textbf{Comparison with GRG.}  GRG's DAG structure enables $O(n \log n + m)$
multiply by memoizing shared haplotype computations.  At 10K diploid, GRG
is 1.3--1.5$\times$ faster than RLE, but at smaller sizes ($\le$2K),
RLE is 2--3$\times$ faster due to lower overhead.  The \texttt{.threads}
format also provides 3$\times$ better compression than GRG.  In practice,
a hybrid strategy could use RLE for moderate sample sizes and GRG for
large-scale biobank analyses.

\textbf{Comparison with RePair.}  Grammar-compressed multiplication
(RePair~\cite{tosoni2022}) is 9--14$\times$ slower than both RLE and GRG at
10K diploid, and its compressed representation is 31$\times$ larger than
\texttt{.threads}.  The grammar decompression overhead dominates the
multiply cost.

For applications requiring many multiply operations (e.g., RHE with
$B = 100$ random vectors at $n = 10{,}000$ diploid), the total compute
for multiplications drops from $\sim$11 seconds (tree) to $\sim$0.6 seconds
(RLE), while preparation time is comparable ($\sim$400\,ms each).

The \texttt{.threads} format remains optimal for storage ($\sim$500$\times$
compression vs dense), with the RLE structure serving as a runtime-only
acceleration layer.  The optional simplification (stripping tmrcas) provides
an additional $\sim$30\% file size reduction for multiply-only workloads.

\begin{thebibliography}{9}
\bibitem{gunnarsson2025}
B.~Gunnarsson, M.~Palacios, P.~P.~Palacios, and P.~F.~Palamara.
\newblock Inferring genome-wide genealogies for large samples.
\newblock \emph{In preparation}, 2025.

\bibitem{dehaas2024}
A.~DeHaas, A.~Biddanda, and D.~Taliun.
\newblock Genotype Representation Graphs for lossless compression of population-scale genetic variation.
\newblock \emph{bioRxiv}, 2024.

\bibitem{pazokitoroudi2020}
A.~Pazoki Toroudi, A.~Chiu, K.~Engleman, et~al.
\newblock Efficient and scalable estimation of heritability with kernel methods.
\newblock \emph{Nature Communications}, 11:3189, 2020.

\bibitem{tosoni2022}
A.~Tosoni, D.~Denti, and S.~Rossi.
\newblock Grammar-compressed matrices for efficient linear algebra.
\newblock \emph{arXiv}, 2022.
\end{thebibliography}

\end{document}
